\section{Bounds for sets without rational geometric progressions}\label{sec:density}
A three-term geometric progression of positive integers has the form
$a<b<c$ with $b^2=ac$; its ratio $b/a$ may be a noninteger. Let
$\overline\alpha$ be the supremum of $\limsup_{N\to\infty}|A\cap[1,N]|/N$
over subsets of the positive integers containing no such triple.
For $s\ge2$ write $S_s$ for the positive integers supported on primes at
most $s$, including $1$. Define $\overline\alpha_s$ by forbidding just
those progressions whose ratio, in lowest terms, has numerator and
denominator in $S_s$. The subscripts below are smoothness cutoffs.

\medskip\noindent\textbf{Theorem.}
\[
  0.7365954 < \overline\alpha < 0.7716011,\qquad
  \overline\alpha_7-\overline\alpha_{11}\ge \frac{64}{139755}.
\]
McNew's published bounds are $0.730027<\overline\alpha<0.772059$
(Math. Comp. 84 (2015), 2893--2910, Theorems 4.3 and 4.5;
\href{https://www.nathanmcnew.com/GPFsets.pdf}{author's text}).
The exact supremum remains undetermined. These claims concern upper
density. The incoming report \texttt{mcnew\_density\_and\_grh\_transfers}
separates the two optima for \emph{natural} density and gives restricted
Rankin optima. A natural-density upper bound in that report cannot be
used as an upper bound for $\overline\alpha$.

The source connection is the replacement of multiplicative patterns by
affine relations in exponents. The pinned readable edition applies
\texttt{exists\_three\_term\_free} and \texttt{exists\_large\_isSparse}
in \texttt{New/Reduction/Data/Existence.lean} to pack places with
controlled multiplicative patterns. There it represents places by
powers and constructs exponents using a Behrend sphere. Here $b^2=ac$
gives the coordinatewise relation $2v_p(b)=v_p(a)+v_p(c)$.
McNew already used smooth-number fibers for this density question.
The additions here are explicit finite templates and a strict,
certified loss when the prime $11$ enters the allowed ratios.
The density proofs below are independent of Dirichlet GRH.

\subsection*{Lower bound by an explicit infinite construction}
The file \texttt{lower.csv} gives thresholds
$1=t_1<\cdots<t_{275}=972405$ and sets $C_i\subset S_7\cap[1,t_i]$.
The checker verifies $|C_i|=i$ and tests $b^2/a\notin C_i$
for all $a<b$ in $C_i$ for which $a\mid b^2$. Thus all $C_i$ are
progression-free. They need not be nested. For $x\ge1$ let
$C(x)$ be $C_i$ for the largest $t_i\le x$, and use the empty set
for $x<1$. The exact sum of the threshold contributions is
\[
 I=\sum_{i=1}^{275}\frac1{t_i}
  =\frac{86690601810929969}{26132271943680000}.
\]

For a rough part $h$ require $(h,210)=1$. At primes $11\le p\le997$
allow exponents $0,1,3,4$ and at all larger primes allow $0,1$.
Both exponent sets contain no nonconstant arithmetic progression of
length three. The allowed $h$ have natural density
\[
 c=\frac{6}{\pi^2}\prod_{p\in\{2,3,5,7\}}\frac{p}{p+1}
          \prod_{11\le p\le997}(1+p^{-3}),
\]
where products indexed by $p$ range over primes. Indeed, start with
squarefree integers, impose exponent zero for $p\le7$, and add exponents
$3,4$ for primes in $[11,997]$. At each such prime the local density
factor is multiplied by $1+p^{-3}$. Finite Euler products give the
density by the Chinese remainder theorem and the tail of the
square-divisibility exceptions has upper density at most
$\sum_{p>Y}p^{-2}$, which tends to zero. Here one can bound the
number of exceptions through $N$ by summing $\lfloor N/p^2\rfloor$.

Take integers $N_0=0$, $N_j=2^{2^j}$ for $j\ge1$.
For each allowed $h$ in $(N_{j-1},N_j]$, insert the numbers
$h C(N_j/h)$ into a single set $A$. The rough part of each inserted
number is $h$. If three inserted numbers formed a geometric progression,
their valuations at each prime greater than $7$ would form an
arithmetic progression in a permitted exponent set, so all would have
the same rough part. Their three smooth parts would lie in the one
progression-free template assigned to that part. This proves avoidance
across all stages as well as within a stage. At $N=N_j$, the cardinality
from the current stage alone is bounded below by
\[
  \sum_{i=1}^{275}\#\{h\le N/t_i:\ h\text{ allowed}\}-275N_{j-1}.
\]
This counts the size of each template via threshold indicators. Since
$N_{j-1}/N_j\to0$, division by $N_j$ gives
$\overline d(A)\ge cI$. Replacing $\pi$ by its strict upper bound
$355/113$ yields the rational lower estimate
\[
 \frac{6\cdot113^2}{355^2}
 \prod_{p\in\{2,3,5,7\}}\frac p{p+1}
 \prod_{11\le p\le997}(1+p^{-3})\,I
   >\frac{7365954}{10^7}.
\]
The checker also certifies $\pi<355/113$ using the alternating series in Machin's formula $\pi=16\arctan(1/5)-4\arctan(1/239)$.

\subsection*{A strict loss when the prime eleven is admitted}
Let $F_s(x)$ be the maximum size of a progression-free subset of
$S_s\cap[1,x]$, and let $g_s(N)$ be the maximum size of a subset of
$[1,N]$ avoiding progressions with $s$-smooth ratio.
McNew's Theorem 4.2 gives $g_s(N)/N\to\overline\alpha_s$.
In a fiber with fixed part $b$ coprime to $2310$, put $x=N/b$.
Allowing all $11$-smooth ratios gives the maximum $F_{11}(x)$ in this
fiber. If only $7$-smooth ratios are forbidden, its maximum is
$H(x)=\sum_{j\ge0}F_7(x/11^j)$ because each $11$-adic layer can be
chosen independently. Always $H(x)\ge F_{11}(x)$.
The following finite comparison gives an additional exclusion throughout
$121\le x<165$.
\[
\begin{array}{c|rrrrrrr}
 [a,b') &[121,125)&[125,132)&[132,135)&[135,140)&[140,154)&[154,162)&[162,165)\\
 \hline
 \max(S_{11}\cap[1,b')) &121&128&132&135&150&160&162\\
 H(a)\ \text{at least} &38&39&40&41&42&43&44\\
 F_{11}(x)\ \text{at most} &37&38&39&40&41&42&43
\end{array}
\]
For each interval, \texttt{intervals.json} supplies three $7$-smooth
sets for layers $j=0,1,2$. Testing their members and their triples
proves the lower entry for $H(a)$. For an upper entry $u$, use the
hypergraph whose vertices are the $11$-smooth numbers up to the listed
maximum, with a hyperedge for every triple $d<e<f$ satisfying $df=e^2$.
A set of $u+1$ surviving vertices would give a cover with budget
$|V|-u-1$ deleted vertices. The seven proof trees in \texttt{trees/}
rule this out by the same exact checks.

At an internal node, branch on a single vertex being deleted or kept.
Discard the covered edges and let $k$ be the remaining deletion budget.
A leaf closes if $k<0$ or a still uncovered edge has all its vertices
fixed to be kept. The remaining type of leaf carries positive integer
weights on uncovered edges. Let $W$ be their sum and $D$ the maximum
total weight incident to any undecided vertex. Every cover using at most
$k$ more deletions covers weight at most $kD$. The supplied weights
satisfy $W>kD$. The verifier checks these integer inequalities, the
two branches at each internal node, all listed members and the
ordering of the smooth numbers. The seven trees contain respectively
$813,3295,1559,1331,2047,3139,2865$ nodes. Their validity is independent
of the numerical optimizer used to find leaf weights.

Fibers with $(b,2310)=1$ partition $[1,N]$ for both optimization
problems. Every $b$ with $N/165<b\le N/121$ costs at least one in the
$11$-smooth problem. Dividing $g_7(N)-g_{11}(N)$ by $N$ and passing
to the limit therefore gives
\[
 \overline\alpha_7-\overline\alpha_{11}
 \ge \frac{\varphi(2310)}{2310}
       \left(\frac1{121}-\frac1{165}\right)=\frac{64}{139755}.
\]
McNew's section 4 obtains the more specific bound
$\overline\alpha_7<0.772059$ immediately before Theorem 4.3.
Consequently
\[
 \overline\alpha\le \overline\alpha_{11}
 < \frac{772059}{10^6}-\frac{64}{139755}
 =\frac{21567021109}{27951000000}<0.7716011.
\]
